% !PRINTED 686
\weberpage{686}
Es soll jetzt bewiesen werden, da\ss{} $\mathfrak{L}$ durch das Verzweigungspolygon $\mathfrak{Z}_\beta$ teilbar ist.

Nehmen wir zun\"achst an, da\ss{} $\mathfrak{A}$ und $\mathfrak{B}$ relativ prim zu $\mathfrak{Z}_\beta$ seien; dann gibt es eine ganze rationale Funktion $x$ von $\beta$, die in keinem Punkt von $\mathfrak{Z}_\beta$ verschwindet, f\"ur die $\omega = x\alpha$ eine ganze Funktion von $\beta$ wird. Setzen wir
\[
f(\omega) = x^n F(\alpha,\beta),
\]
so ist:
\[
f'(\omega) = x^{n-1}F'(\alpha).
\]

Ist $\tau$ die Basisform von $\Omega$ in bezug auf $\beta$, so ist nach \S\ 182 $f'(\omega)$ durch $f'(\tau)$ teilbar, und da nach Voraussetzung $\beta$ in keinem Verzweigungspunkt unendlich wird, so ist $f'(\tau)$ und folglich auch $F'(\alpha)$ durch das Verzweigungspolygon $\mathfrak{Z}_\beta$ teilbar, und unsere Behauptung erwiesen. Wir setzen:
\[
\mathfrak{L} = \mathfrak{R}\mathfrak{Z}_\beta,
\]
\weq{8}{F'(\alpha) = \dfrac{\mathfrak{R}\mathfrak{Z}_\beta}{\mathfrak{A}^{n-2}\mathfrak{B}^m}.}

Diese Formel ist hierdurch nur unter der Voraussetzung erwiesen, da\ss{} weder $\alpha$ noch $\beta$ in einem der Verzweigungspunkte unendlich wird.

Machen wir aber die lineare Substitution:
\weq{9}{\begin{aligned}
\alpha(\alpha_1+1) &= a\alpha_1,\\
\beta(\beta_1+1) &= b\beta_1,
\end{aligned}}
so ist nach \S\ 190
\[
\mathfrak{Z}_\beta = \mathfrak{Z}_{\beta_1},
\]
und wir k\"onnen die Konstanten $a$, $b$ so w\"ahlen, da\ss{} $\alpha_1$ und $\beta_1$ in keinem Punkte von $\mathfrak{Z}_\beta$ unendlich werden. Wir haben nur f\"ur $a$ und $b$ irgend konstante Werte zu nehmen, die $\alpha$ und $\beta$ in keinem Punkte von $\mathfrak{Z}_\beta$ annehmen.

Ist $F_1(\alpha_1,\beta_1)=0$ die zwischen $\alpha_1$, $\beta_1$ bestehende rationale Gleichung, so ist:
\[
F_1(\alpha_1,\beta_1) = (\alpha_1+1)^n(\beta_1+1)^m F(\alpha,\beta),
\]
oder auch nach (9):
\[
\alpha^n\beta^m F_1(\alpha_1,\beta_1) = a^n b^m \alpha_1^n \beta_1^m F(\alpha,\beta);
\]
ferner $d\alpha : d\alpha_1 = \alpha^2 : a\alpha_1^2$, und folglich, mit R\"ucksicht auf $F=0$, $F_1=0$:
\weq{10}{\alpha^{n-2}\beta^m F_1'(\alpha_1) = a^{n-1}b^m \alpha_1^{n-2}\beta_1^m F'(\alpha).}
